\begin{abstract}
We establish two structural foundations for the deployment-safety
theorem of \citep[Deployment Safety]{lasser2026paper10} in the
Goal-Frontier Maximization sequence. First, bounded co-evolution
is derived as a corollary of clipped-LLR SPRT machinery
((C5.HOEFF) of \citep[Deployment Safety]{lasser2026paper10}), a
new upper exposure-rate cap (C7.RATE), and the verification
protocol's channel-projection structure: the per-channel coupling
magnitude $\bar{M}^{\mathrm{cum}}$ is bounded by a closed-form
constant $C \cdot \Bclip \cdot \lammax \cdot \taumeta$ with
$C \in \{1, \Kch - 1\}$, all factors deployment-class quantities
intensive in $|P|$. Second, a scoped Concentration-Gap selection
theorem proves that under three operationally auditable structural
conditions --- representativeness (REP), bounded dispersion (DISP),
coalition closure (COAL) --- and a Lipschitz embedding
compatibility assumption, the proxy-truth Goodhart slack is
bounded by an explicit $\chi^2$-divergence quantity that controls
the Herfindahl-index trade-flow concentration through
$\chi^2(w \,\|\, \mu) = N \HHI(w) - 1$ for uniform base measure
$\mu$. The selection theorem replaces the monolithic HHI
surrogate-adequacy assumption of
\citep[Deployment Safety]{lasser2026paper10} with three concrete
checkable conditions, and converts the Concentration-Gap
conjecture inherited from
\citep[Microfoundation]{lasser2026micro} into a scoped theorem
with named structural prerequisites. Audit procedures for
(C7.RATE), (REP), (DISP), (COAL), and the embedding compatibility
precondition are specified.
\end{abstract}
